% Stylized illustration of cosine similarity between two document vectors (greyscale).
% Intended to be \input{} into a document that loads tikz with arrows.meta, calc.
\begin{tikzpicture}[
    >=Stealth,
    docvec/.style = {-{Stealth[length=7pt]}, line width=1.1pt, black},
    lab/.style    = {font=\footnotesize, align=left}
  ]

  % shared bounding box so this panel aligns with kingqueen.tex
  \useasboundingbox (-0.3,-1.0) rectangle (4.9,3.7);

  \coordinate (O) at (0,0);

  % --- faint embedding-space axes ----------------------------------------
  \draw[gray!50, -{Stealth[length=5pt]}] (O) -- (4.1,0);
  \draw[gray!50, -{Stealth[length=5pt]}] (O) -- (0,3.7);

  % --- two document vectors (solid vs dashed to tell them apart) ----------
  \coordinate (p1) at (30:3.4);
  \coordinate (p2) at (62:3.0);
  \draw[docvec]         (O) -- (p1);
  \draw[docvec, dashed] (O) -- (p2);

  % --- angle between them -------------------------------------------------
  \draw[black, line width=0.8pt] (30:1.05) arc (30:62:1.05);
  \node[font=\small] at (46:1.45) {$\theta$};

  % --- labels at the vector tips -----------------------------------------
  \node[lab, anchor=west]       at (p1) {Scientific\\paper 1};
  \node[lab, anchor=south west] at (p2) {Scientific\\paper 2};

  % --- annotation ---------------------------------------------------------
  \node[anchor=west, font=\small] at (-0.2,-0.6)
        {$\cos\theta$: similarity of two papers};

\end{tikzpicture}
